\anexoItem{Especificación detallada de requerimientos no funcionales}{
    \begingroup
    \renewcommand{\baselinestretch}{1}
    \footnotesize
    \setlength{\tabcolsep}{4pt}
    \renewcommand{\arraystretch}{1.2}

    \begin{longtable}{@{}p{1.4cm}p{4.4cm}p{8.2cm}c@{}}
        \caption[]{Especificación detallada de requerimientos no funcionales}%
        \addcontentsline{lot}{table}{Tabla~\thetable. Especificación detallada de requerimientos no funcionales}%
        \label{tab:anexo_especificacion_rnf} \\
        \toprule
        \textbf{Código} & \textbf{Nombre del Requerimiento} & \textbf{Descripción del Requerimiento No Funcional} & \textbf{Prioridad} \\ \midrule
        \endfirsthead
        
        \toprule
        \textbf{Código} & \textbf{Nombre del Requerimiento} & \textbf{Descripción del Requerimiento No Funcional} & \textbf{Prioridad} \\ \midrule
        \endhead
        
        \midrule
        \multicolumn{4}{r}{\textit{(Continúa en la siguiente página)}} \\
        \endfoot
        
        \bottomrule
        \multicolumn{4}{@{}p{\linewidth}@{}}{\vspace{3pt}\textit{Nota.} Criterios de calidad, seguridad, rendimiento y arquitectura de la solución tecnológica. Elaboración propia.}
        \endlastfoot

        RNF-01 & 
        Optimizar el tiempo de respuesta del motor de cálculo & 
        Procesar la nómina completa y generar los snapshots de resultados para un volumen de hasta 1,000 colaboradores en un tiempo inferior a 10 segundos utilizando Worker Threads (Piscina). &
        Alta \\ \addlinespace

        RNF-02 & 
        Minimizar la latencia en endpoints transaccionales HTTP & 
        Responder a las solicitudes HTTP de consulta y mutaciones ligeras a través del motor Fastify 5 con una latencia media menor a 300 milisegundos en condiciones operativas estándar. &
        Alta \\ \addlinespace

        RNF-03 & 
        Procesar tareas intensivas en segundo plano de forma asíncrona & 
        Delegar operaciones pesadas (cálculo de nómina, validación de archivos masivos de más de 5,000 filas y renderizado de boletas) a colas en memoria con BullMQ y Redis 7 sin bloquear el Event Loop de Node.js. &
        Alta \\ \addlinespace

        RNF-04 & 
        Cifrar datos financieros sensibles en reposo & 
        Encriptar simétricamente los números de cuenta bancaria sueldo, códigos CCI y cuentas CTS en la base de datos utilizando el algoritmo AES-256-GCM en la capa de aplicación. &
        Alta \\ \addlinespace

        RNF-05 & 
        Enmascarar información financiera en interfaces y respuestas & 
        Ocultar dinámicamente los dígitos intermedios de las cuentas bancarias en las respuestas JSON y componentes de vista, mostrando únicamente los últimos 4 dígitos visibles para verificación del usuario. &
        Media \\ \addlinespace

        RNF-06 & 
        Proteger boletas PDF mediante cifrado y firma digital & 
        Cifrar individualmente las boletas de pago generadas en formato PDF utilizando el número de DNI del trabajador como clave de apertura e incorporar firma digital institucional. &
        Alta \\ \addlinespace

        RNF-07 & 
        Garantizar la conservación de parámetros históricos de ley & 
        Preservar el registro histórico de comisiones de AFP, valores de UIT y RMV bajo el esquema de Slowly Changing Dimensions (SCD Tipo 2), permitiendo auditar periodos pasados sin colisión con los valores actuales. &
        Alta \\ \addlinespace

        RNF-08 & 
        Mantener alta disponibilidad durante los cierres mensuales & 
        Garantizar una disponibilidad del servicio (Uptime) de al menos 99.5\,\% durante los días hábiles del periodo de cierre de nóminas y pagos bancarios. &
        Alta \\ \addlinespace

        RNF-09 & 
        Asegurar la resiliencia ante caídas de servicios externos & 
        Implementar límites de espera (Timeout de 5 segundos) y manejo de excepciones ante fallos de la API de RENIEC o webhooks de n8n, garantizando que el sistema no interrumpa su operación. &
        Media \\
    \end{longtable}
    \endgroup
}